% =====================================================================
% Cover note for paper/kraken.pdf.
%
% WHY THIS FILE EXISTS. The report is 130 pages. That length is defensible --
% it is four development cycles with the failed arms left in -- but it is the
% wrong first object to put in front of someone who has not asked for it. This
% is the two pages that go in the email body's place, and it leads with the
% retraction rather than the result, because the retraction is the better
% evidence about how the work was done.
%
% Build:  pdflatex -interaction=nonstopmode cover_note.tex   (one pass; no
% refs, no citations, nothing to settle).
% =====================================================================
\documentclass[11pt,a4paper]{article}

\usepackage[T1]{fontenc}
\usepackage{mathptmx}
\usepackage[scaled=0.92]{helvet}
\usepackage[margin=1in,top=1in,bottom=1in]{geometry}
\usepackage{microtype}
\usepackage{booktabs}
\usepackage[colorlinks=true,linkcolor=blue,urlcolor=blue]{hyperref}
\usepackage{titlesec}
\usepackage{enumitem}

\titlespacing*{\section}{0pt}{1.9ex plus .5ex minus .2ex}{0.9ex plus .2ex}
\titleformat{\section}{\normalfont\bfseries\normalsize}{}{0pt}{}
\setlist[itemize]{itemsep=2pt,topsep=4pt,parsep=0pt,leftmargin=1.3em}
\pagestyle{empty}
\setlength{\parskip}{0.55em}
\setlength{\parindent}{0pt}

\begin{document}

{\bfseries\large Exact Inference Is Not Enough}\\[0.15em]
{\normalsize Opponent models and endgames in six-player Literature ---
KRAKEN v1.1}\\[0.35em]
{\small Kausthubh Veldanda \quad$\cdot$\quad accompanying
\texttt{kraken.pdf}, 130 pages}

\vspace{0.4em}
\hrule
\vspace{0.2em}

Literature (also Fish, or Canadian Fish) is a six-player, two-team,
imperfect-information card game. Players interrogate opponents for named cards
and score by declaring exactly how a six-card half-suit is split among their
own team. Every action is public, so the whole hidden state is the deal --- and
there is no legal communication channel, which puts the game outside both the
two-player zero-sum machinery that solved poker and the fully cooperative
machinery built for Hanabi. This report is four development cycles on an engine
for it, written with the arms that failed left in.

I am writing because I would value criticism of the method, and because one
episode in it seems worth more than the engine is.

\section{The result I retracted, and how}

An earlier version of this work reported that KRAKEN beat an independently
developed C++ engine by $+2.347$ sets per game over $10{,}000$ audited games.
\textbf{That figure is withdrawn.} It measured my bridge, not my engine.

The bridge was stateless: it replayed the entire public log into a freshly
reset copy of their agent at every decision, where their own arbiter keeps one
agent per seat alive for the whole deal. Those are the same thing only if their
agent is a pure function of (reset state, event sequence). It is not.

What made this findable was not suspicion of the number. It was a control. To
check that their engine had not simply got weaker, I ran their whole released
ladder, v0.2 through v0.7, and the ladder came out \emph{backwards}: in my
arbiter their later releases scored worse, and their declaration error climbed
monotonically with their own version number (Spearman $-1.000$ against $+0.771$
the other way). A project's successive releases getting steadily worse at the
one thing the game scores is not a shape that improving strength produces. So
the ladder was the instrument, not the finding.

Localising it took a decision-level parity test --- the same game put to both
transports, decision by decision:

\begin{center}
\small
\begin{tabular}{lrrr}
\toprule
their version & decisions & divergent & rate \\
\midrule
v0.2 (scripted) & $323$ & $\mathbf{0}$ & $0.00\%$ \\
v0.3 (scripted) & $338$ & $\mathbf{0}$ & $0.00\%$ \\
v0.4 & $415$ & $81$ & $19.52\%$ \\
v0.5 & $336$ & $137$ & $40.77\%$ \\
v0.6 & $292$ & $36$ & $12.33\%$ \\
v0.7 & $296$ & $84$ & $28.38\%$ \\
\bottomrule
\end{tabular}
\end{center}

Exactly zero on both of their scripted baselines, and large on every release
from v0.4 --- the version their fitted belief arrives at. A clean boundary, and
the same rung at which the two arbiters' ladders stop agreeing.

Then it was priced, under a paired design that holds the deal, the seats, the
agent seeds and my own engine fixed and changes only whether their process
survives between decisions. \textbf{The bridge was worth $+3.18$ $[+3.01,
+3.34]$ sets a game to me} --- more than the entire margin it produced. Their
declaration errors fall roughly tenfold, from $0.880$ a game to $0.092$, once
their engine is allowed to keep the state its own arbiter lets it keep.

Corrected, KRAKEN \textbf{loses} to that engine by $-0.70$. Measured
independently inside \emph{their} arbiter, through a bot-package protocol they
published, by $-0.62$. Two routes that share no host, no rules implementation,
no deal generator and no line of code agree to within $0.08$ --- where the
published number was out by three and had the sign wrong.

The general lesson is one the report had already written down in an appendix
and then not applied to its own headline: \emph{an absolute measured through a
bridge is a statement about the bridge as well as about the engines.} Paired
contrasts that hold the bridge fixed on both sides were unaffected, and are
what the rest of the report rests on.

\section{The central result is also negative}

Hidden-state inference here reduces to counting assignments in a
degree-constrained bipartite system. That is $\#P$-complete in general, but
real positions carry only $4.1$ distinct candidate sets across $25.1$ unlocated
cards, so a dynamic program returns the exact partition function, marginals,
uniform samples and joint probabilities --- validated against exhaustive
enumeration.

Making the posterior exactly correct \textbf{changed nothing in play}
($-0.008$ sets per duplicate deal-pair over $250$ pairs) and made it a
\emph{worse} predictor of where cards actually sit than the biased sampler it
replaced ($1.3618$ against $1.3408$ mean negative log-likelihood). A uniform
posterior is a correct theorem about a false hypothesis: it assumes choices
carry no information beyond legality. The old sampler's bias had accidentally
encoded an opponent model, and removing the bias removed the model. Replacing
that accident with an explicit one-parameter choice model is worth $+1.92$
$[+1.48, +2.36]$ sets per pair --- the largest single gain in the study.

Two of the results that follow. Residual errors are not failures to \emph{find} cards:
$95.3\%$ are \emph{allocation} errors, where the team held all six of a
half-suit and split them wrongly. Handing a seat the true deal one side at a
time --- a ceiling obtained by cheating, never a strength figure --- prices
teammates' cards at $+3.41$ against opponents' $+1.31$, and I had predicted the
reverse.

\section{How the claims are held to account}

Forty-three pre-registrations, written before the runs, each naming its
prediction and its ship bar; Appendix~A tabulates every one on the same
footing, \emph{including} those whose predictions failed, and a test checks
that table against the directory both ways. Nothing ships below $\pm0.15$ sets
per game. Every number in the report is pinned by a script to the artifact
that produced it, and the script is verified by perturbing an artifact and
confirming it fails. Cross-engine work is pinned to
an upstream commit and re-pinned only on evidence: $6{,}329$ decisions replayed
through binaries built at both ends, zero mismatches.

None of that caught the bridge. What caught it was running a control that had
no reason to be interesting.

\section{What is open}

Which piece of their state the bridge discards is only partly separated, and
what I have went against the convenient reading. Their determinized search is
the largest term I can name, and the one a fix could reach --- their reset
re-seeds it. But sorting the divergences into disagreements about \emph{which}
ask to make and disagreements about \emph{whether to declare at all} puts
everything I can identify in the first, while the second --- the band that
costs sets, since a declaration taken a turn early is a wrong one --- does not
move at all. Two mechanisms I had believed in did not survive their source, and
$13.92\%$ of decisions still differ with everything I found switched off. So
the bridge cannot be repaired by re-seeding, and I do not yet know what the
remainder is. The $+3.18$ price rests on $2{,}000$ pairings: an order of
magnitude past its interval, but I would want more games before quoting it
alone.

I would be glad of any reaction, and particularly of the sharpest version of
``that control only worked by luck.'' I think it did, and I do not yet know
what to put in its place.

\vspace{0.4em}
\hrule
\vspace{0.1em}

{\small Report, code, pre-registrations and result artifacts:
\url{https://github.com/kv1514/fish-researchp12}}

\end{document}
